\subsection{Adaptation on an external source and the remaining validation gap}
\label{sec:code-discussion}
The completed CODE analysis provides a matched patient-held-out comparison within CODE after supervised adaptation, which the historical selected profiles do not supply. The gain is conditional on supervised adaptation, length-aware preprocessing, a new head and calibration-selected operating points. Because these components changed together, the comparison does not isolate the causal contribution of Transformer attention, morphology or one preprocessing choice. Frozen P6/P8 controls provide reference systems evaluated on the same ECGs; they do not constitute a comparison with a matched adaptation budget.

The five provider flags limit the scope of the experiment. It includes no PVC evaluation, full P6/P8 profile assessment, P4 assessment, new expert adjudication or clinical outcome measurement. The candidate showed sensitivity losses for AFIB and SB and lower 1AVB PPV relative to P8, despite the improved average. The failed promotion decision is therefore retained. Better mean Brier/ECE and high mean binary accuracy cannot replace endpoint-specific error review or a clinical utility study.

The independent implementation readback supports numerical consistency for one sealed prediction artifact; it does not add patients or institutions. Likewise, zero within-CODE patient-key overlap and exact waveform screening do not certify all historical or cross-database separation. The earlier PR~6 CODE screening involved a different cohort (24,682 ECGs, 17,300 patients), preprocessing and candidate history. Its scores must not be pooled or paired with scores from this 4,921-ECG experiment, or subtracted from them as though the observations were interchangeable. The studies document distinct stages. An independent institutional cohort with a frozen model, complete overlap accounting and compatible reference labels remains necessary to assess transport of the adapted system.
